% titel: Blatt 6, A5: Eulersche Formel
% Teil II: der durchgerechnete Weg.

\subsection*{Gegeben und gesucht}
Mit $e^{it}=\cos t+i\sin t$ zeigen: (i) $|e^{it}|=1$, (ii) $e^{-it}$ ist
das Spiegelbild von $e^{it}$ an der $x$-Achse, (iii) $(e^{it})^n=e^{int}$
(Induktion). Danach aus $n=2$ die Doppelwinkelformeln folgern.

\subsection*{a) Eigenschaften}
\begin{itemize}
  \item[(i)] \emph{Gedankenschritt:} Betrag aus $|z|^2=z\bar z$ und der
        Kreisfunktionsidentität. Wegen $\overline{e^{it}}=\cos t-i\sin t$:
        \begin{equation*}
          |e^{it}|^2=e^{it}\overline{e^{it}}
          =(\cos t+i\sin t)(\cos t-i\sin t)
          =\cos^2t+\sin^2t=1,
        \end{equation*}
        also $|e^{it}|=1$. (Man kann auch $e^{it}e^{-it}=e^0=1$ nutzen.)
  \item[(ii)] \emph{Gedankenschritt:} Konjugation ersetzt $i$ durch $-i$:
        \begin{equation*}
          \overline{e^{it}}=\cos t-i\sin t=e^{-it} .
        \end{equation*}
        Geometrisch: Der Punkt zum Winkel $t$ gespiegelt an der $x$-Achse
        ist der Punkt zum Winkel $-t$.
  \item[(iii)] \emph{Gedankenschritt:} Induktion über $n$, mit der
        Additionsregel $e^{a}e^{b}=e^{a+b}$. Induktionsanfang $n=1$
        trivial; Induktionsschritt $n\to n+1$:
        \begin{equation*}
          (e^{it})^{n+1}=(e^{it})^n\,e^{it}=e^{int}\,e^{it}
          =e^{i(n+1)t}.
        \end{equation*}
        Damit gilt $(e^{it})^n=e^{int}$ für alle $n\in\mathbb N$.
\end{itemize}

\subsection*{b) Doppelwinkelformeln aus $n=2$}
\emph{Gedankenschritt:} Zwei Darstellungen derselben Zahl gleichsetzen und
Real-/Imaginärteil getrennt lesen. Aus dem Potenzgesetz $(e^{it})^2=e^{i2t}$:
\begin{align*}
  (\cos t+i\sin t)^2
  &=\cos^2t+2i\sin t\cos t-\sin^2t
  =\cos^2t-\sin^2t+2i\sin t\cos t,\\
  e^{i2t}&=\cos 2t+i\sin 2t .
\end{align*}
Vergleich der Realteile: $\cos 2t=\cos^2t-\sin^2t$. Vergleich der
Imaginärteile:
\begin{equation*}
  \boxed{\cos 2t=\cos^2t-\sin^2t},\qquad
  \boxed{\sin 2t=2\sin t\cos t}.
\end{equation*}

\begin{ergebnis}
\begin{equation*}
  |e^{it}|=1,\quad \overline{e^{it}}=e^{-it},\quad (e^{it})^n=e^{int},\quad
  \cos2t=\cos^2t-\sin^2t,\quad \sin2t=2\sin t\cos t .
\end{equation*}
\end{ergebnis}

\begin{kontrolle}
\textbf{Beispieldaten:} $t=\frac\pi4$: $\sin\frac\pi2=1$, rechts
$2\sin\frac\pi4\cos\frac\pi4=2\cdot\frac{\sqrt2}{2}\cdot\frac{\sqrt2}{2}=1$.
\checkmark $t=\frac\pi2$: $\cos\pi=-1$, rechts
$\cos^2\frac\pi2-\sin^2\frac\pi2=0-1=-1$. \checkmark

\textbf{Grenzfall:} $|e^{it}|=1$ für alle reellen $t$ -- selbst für
riesiges $t$ bleibt $e^{it}$ auf dem Einheitskreis (kein Wachstum über
$1$). Das illustriert den Unterschied zur reellen Funktion $e^x$, die
stark wächst. \checkmark
\end{kontrolle}

\begin{merksatz}
Die Eulersche Formel verbindet Potenzieren und Winkeldrehen: $e^{it}$ hat
Betrag $1$ und liegt auf dem Einheitskreis; Konjugation spiegelt an der
reellen Achse und liefert $e^{-it}$; $(e^{it})^n=e^{int}$. Setzt man das
Quadrat in Real-/Imaginärteile um, erhält man die Doppelwinkelformeln.
\end{merksatz}
